% WIP draft for Section 4.1 (rigidity first proof).
\section{A Classical Rigidity Corollary}\label{sec:classical-rigidity}

The localization operator $\mathcal{L}_U$ on the Bargmann--Fock space is
diagonalized by Hermite functions whenever $U$ is a disk centered at the
origin: this is an immediate consequence of the rotational symmetry of $U$
together with the radial structure of the Hermite weight. The converse
question, raised by Abreu and D\"orfler, asks whether eigenfunction status of
even a \emph{single} Hermite function already forces the underlying domain to
be a disk. The next theorem records a positive answer in the simply connected
$C^2$ regime.

\begin{theorem}[Abreu--D\"orfler rigidity]\label{thm:eig-loc-U}
Let $U\subset \C$ be a simply connected bounded domain with $C^2$ boundary,
and let $h_k$ be a Hermite function ($k\ge 0$). If $h_k$ is an eigenfunction
of the localization operator $\mathcal{L}_U$, then $U$ is a disk.
\end{theorem}

We will give two proofs of Theorem~\ref{thm:eig-loc-U}. The first, presented
below, leverages the Fourier-analytic machinery developed earlier in the
paper, reducing the question to a free boundary problem for a logarithmic
potential. The second, in the next subsection, proceeds along more
classical complex-analytic lines.

\subsection{First proof}

\begin{proof}[First proof under the additional assumption that $\partial U$ is $C^2$]
The Bargmann transform of the Hermite function $h_k$ is a nonzero scalar
multiple of the monomial $z^k$. Thus, arguing as in
\cite{Abreu-Doerfler}, the assumption that $h_k$ is an eigenfunction of
the localization operator $\mathcal{L}_U$ is equivalent to the identity
\[
\int_U z^k\, e^{\pi w\overline{z}}\, e^{-\pi |z|^2}\,dA(z)
=\lambda\, w^k,
\qquad w\in \C.
\]
Both sides are entire functions of $w$. Since $U$ is bounded and the
weight $e^{-\pi|z|^2}$ is integrable, the dominated convergence
theorem applied to difference quotients shows that for each
multi-index $\alpha$ in $w$, the partial derivative
$\partial_w^\alpha$ commutes with the integral on the left-hand side;
the absolute integrability of $|z|^{|\alpha|+k} e^{\pi
|w||z|}e^{-\pi|z|^2}\mathbf{1}_U$ for fixed $w$ provides the required
local majorant. Iterating $k$ times, we obtain
\[
\frac{1}{(\pi)^k}\,\partial_w^k
\Big[\int_U z^k\, e^{\pi w\overline{z}}\, e^{-\pi |z|^2}\,dA(z)\Big]
=\int_U z^k\overline{z}^k\, e^{\pi w\overline{z}}\, e^{-\pi |z|^2}\,dA(z)
=\int_U |z|^{2k}\, e^{\pi w\overline{z}}\, e^{-\pi |z|^2}\,dA(z).
\]
On the right-hand side, $\partial_w^k(\lambda w^k)=\lambda\,k!$.
Setting $c(\lambda,k):=\lambda\,k!/\pi^k\in\R$ and using
$\mathbf{1}_U$ to extend the integration to $\C$, we obtain
\begin{equation}\label{eq:fourier-delta}
\int_{\C} \mathbf{1}_U(z)\,|z|^{2k}\,e^{-\pi |z|^2}\, e^{\pi \overline{z}\, w}\,dA(z)
=c(\lambda,k),
\end{equation}
for some real constant $c(\lambda,k)$.

Consider now the compactly supported distribution
\[
\eta:=|z|^{2k}e^{-\pi |z|^2}\mathbf{1}_U(z)-c(\lambda,k)\,\delta_0,
\]
which has compact support in any closed ball containing $U\cup\{0\}$.
Writing $z=x+iy$ and $w\in \C$, we may rewrite the exponential factor in
\eqref{eq:fourier-delta} as
\[
e^{\pi \overline{z}\, w}=e^{\pi(x-iy)w}=e^{\pi xw}\, e^{-i\pi yw}.
\]
Hence \eqref{eq:fourier-delta} implies that the Fourier transform
$\widehat{\eta}$, viewed as an entire function on $\C^2$ via Paley--Wiener,
satisfies
\[
\widehat{\eta}(w,-iw)=0,\qquad \forall w\in\C.
\]
Since $c(\lambda,k)\in \R$, taking complex conjugates in
\eqref{eq:fourier-delta} after replacing $w$ by $\overline{w}$ also yields
\[
\widehat{\eta}(w,iw)=0,\qquad \forall w\in\C.
\]

Applying Lemma~\ref{lemma:fourier-to-fboundary} to the distribution
$\eta$, we obtain a function $u\in L^2(\R^2)$ such that
\begin{equation}\label{eq:potato}
\Delta u=f(|z|)\,\mathbf{1}_U-c(\lambda,k)\,\delta_0
\end{equation}
in the sense of distributions, where
\[
f(|z|)=|z|^{2k}e^{-\pi |z|^2}.
\]
Writing $u$ explicitly as a logarithmic potential of the right-hand side
of \eqref{eq:potato}, we conclude that $U$ satisfies
\begin{equation}\label{eq:potential}
\int_U \log|z-z'|\, f(|z'|)\,dA(z')
=c(\lambda,k)\,\log|z|,
\qquad z\in \C\setminus (U\cup\{0\}).
\end{equation}

\medskip

\noindent\emph{Reduction to $0\notin U$.}\quad
Suppose first that $0\in U$, and pick $\rho>0$ small enough that
$\overline{D_\rho(0)}\subset U$. The radial measure
$f(|z|)\,\mathbf{1}_{D_\rho(0)}\,dA$ generates a logarithmic potential
\[
\Phi_\rho(z):=\int_{D_\rho(0)} \log|z-z'|\, f(|z'|)\,dA(z'),
\]
which, by Newton's shell theorem in the plane, is radial; outside
$D_\rho(0)$ it equals
\[
\Phi_\rho(z)=\Big(\int_{D_\rho(0)}f(|z'|)\,dA(z')\Big)\log|z|+\text{const}.
\]
Subtracting $D_\rho(0)$ from $U$ thus preserves identity
\eqref{eq:potential} on the complement of the new set, at the cost of
replacing the constant $c(\lambda,k)$ by another real constant. Hence we
may assume $0\notin U$. Moreover, $0\notin \partial U$, because the
left-hand side of \eqref{eq:potential} stays bounded as $z\to 0$, whereas
the right-hand side diverges to $-\infty$ if $c(\lambda,k)\neq 0$.
The case $c(\lambda,k)=0$ is incompatible with the integral being
positive at $z=0$, so it is also excluded.

\medskip

Define the logarithmic potential
\[
\Psi_U(z):=\int_U \log|z-z'|\, f(|z'|)\,dA(z').
\]
We claim that $\Psi_U$ is \emph{locally radial}: for each $z_0\in U$,
there exists a ball $B(z_0)\subset U$ such that $\Psi_U$ agrees on
$B(z_0)$ with a radial function of $|z|$.

To prove this, let $B\supset U$ be a ball centered at $0$, and define
\[
\Phi(z):=\int_B \log|z-z'|\, f(|z'|)\,dA(z').
\]
Since both $B$ and the weight $f(|z'|)$ are radial, $\Phi$ is radial; in
particular, $\nabla \Phi(z)$ is parallel to $z$ for every $z\in \R^2$.
Furthermore, on the interior of $B$,
\[
\Delta \Phi(z)=2\pi\, f(|z|).
\]

Now $\Psi_U$ satisfies \eqref{eq:potential} on the unbounded
component of $\C\setminus(U\cup\{0\})$. Differentiating that identity
in $z$ for $z\not\in \overline{U}\cup\{0\}$ gives
\[
\nabla \Psi_U(z)=c(\lambda,k)\,\frac{z}{|z|^2},
\qquad z\in \C\setminus(\overline{U}\cup\{0\}).
\]
The classical jump relations for the single-layer / volume logarithmic
potential of an $L^\infty$ density on a $C^2$ domain (see, e.g.,
\cite{Gustafsson}) imply that $\nabla\Psi_U$ extends continuously
across $\partial U$ from \emph{both} sides, and the boundary values
agree. Hence on $\partial U$,
\[
\nabla \Psi_U(z)=c(\lambda,k)\,\frac{z}{|z|^2},
\qquad z\in \partial U,
\]
which is, in particular, parallel to $z$. Therefore the difference
\[
v:=\Psi_U-\Phi
\]
satisfies $\Delta v=2\pi f(|z|)\mathbf 1_U-2\pi f(|z|)\mathbf 1_B
=-2\pi f(|z|)\mathbf 1_{B\setminus U}$ in the sense of distributions
on $B$, so $v$ is harmonic in $U$ and in $B\setminus\overline{U}$, and
on $\partial U$ satisfies
\[
\nabla v(x)\ \text{is parallel to}\ x.
\]
The next lemma converts this gradient condition into local radiality of $v$.

\begin{lemma}\label{lem:harmonic-locally-radial}
Let $V\subset \R^2$ be a domain, and suppose there exists a harmonic
function $v\in C^1(\overline{V})$ such that $\nabla v(x)$ is parallel to
$x$ for every $x\in \partial V$. Then $\nabla v(x)$ is parallel to $x$
for every $x\in V$. In particular, on every circle
$S=\{|x|=r\}\subset \R^2$, the function $v$ is constant on every
connected component of $S\cap V$.
\end{lemma}

\begin{proof}
Consider the angular derivative
\[
g(x):=x_1\,\partial_2 v(x)-x_2\,\partial_1 v(x)
=\partial_\theta v(x),
\]
which equals the tangential derivative of $v$ along the circle
through $x$ centered at $0$. A direct computation gives
\[
\Delta g
=x_1\,\Delta(\partial_2 v)-x_2\,\Delta(\partial_1 v)
+2\big(\partial_1\partial_2 v-\partial_2\partial_1 v\big)
=0,
\]
since $v$ is harmonic and partial derivatives of harmonic functions
are themselves harmonic. By the boundary hypothesis, $\nabla v(x)$ is
parallel to $x$ on $\partial V$, so $g\equiv 0$ on $\partial V$.
Together with $g\in C(\overline{V})$, the maximum principle for
harmonic functions yields $g\equiv 0$ in $V$. Hence $\nabla v(x)$ is
parallel to $x$ at every interior point as well, and on any circle
$S=\{|x|=r\}$ the tangential derivative $\partial_\theta v$ vanishes
on $S\cap V$. Therefore $v$ is constant on each connected component
of $S\cap V$.
\end{proof}

Applying Lemma~\ref{lem:harmonic-locally-radial} on small subdomains
$V\subset U$ where $v$ is harmonic and $C^1$ up to the boundary, we
conclude that $v$ is locally radial on $U$. Since $\Phi$ is globally
radial, so is $\Psi_U=\Phi+v$, locally on $U$.

\begin{lemma}\label{lem:local-to-global-radial}
Let \(W\subset \R^2\setminus\{0\}\) be a connected open set, and let
\(u\in C(W)\). Suppose that for every \(x\in W\) there exist \(r_x>0\)
and a scalar function \(\phi_x\) such that
\[
u(y)=\phi_x(|y|)
\qquad\text{for all } y\in B(x,r_x)\subset W.
\]
Then there exists a single scalar function \(\phi\), defined on the
interval
\[
I_W:=\{|y|:y\in W\},
\]
such that
\[
u(y)=\phi(|y|)
\qquad\text{for all } y\in W.
\]
\end{lemma}

\begin{proof}
For each \(x\in W\), choose a ball \(B_x:=B(x,r_x)\subset W\) on which
\(u\) is a radial function of \(|y|\), with profile $\phi_x$. We
first note that whenever $B_x\cap B_{x'}\ne\emptyset$,
\[
\phi_x(|y|)=u(y)=\phi_{x'}(|y|)
\qquad\text{for every }y\in B_x\cap B_{x'}.
\]
Since $B_x\cap B_{x'}$ is open and avoids the origin, the map $y\mapsto
|y|$ is a submersion on it, so the set of radii it attains contains a
nontrivial open interval on which $\phi_x$ and $\phi_{x'}$ agree.

Fix $x_*\in W$ and define a candidate function $\phi$ on $I_W$ as
follows. Given $y\in W$, connect $x_*$ to $y$ by a polygonal path in
$W$, which exists because $W$ is open and connected. Its image is
compact, hence covered by finitely many balls
$B_{x_1},\dots,B_{x_m}$ with successive nonempty overlaps and
$x_1=x_*$, $y\in B_{x_m}$. The previous paragraph implies $\phi_{x_j}$
and $\phi_{x_{j+1}}$ agree on a nontrivial interval of radii, so by
analytic-continuation along the chain we may unambiguously set
$\phi(|y|):=\phi_{x_m}(|y|)$. Independence of the chain follows from
the same overlap argument: any two chains from $x_*$ to $y$ can be
interleaved into a single chain, and successive overlaps force the
final value to coincide. Hence $\phi$ is well-defined on $I_W$ and
satisfies $u(y)=\phi(|y|)$ for every $y\in W$.
\end{proof}

We now conclude the proof. Since \(U\) is simply connected and bounded,
it suffices to rule out the possibility that \(U\) fails to be a disk
centered at the origin. Suppose therefore, for contradiction, that
\(U\) is not such a disk. Then \(\partial U\) is not a single circle
centered at $0$, so there exist radii \(0<r_1<r_2\) such that the
annulus
\[
A:=\{z\in \R^2:\ r_1<|z|<r_2\}
\]
meets both $U$ and $\partial U$ in a non-trivial way; in particular,
we may pick a point
\[
x_0\in \partial U\cap A.
\]
Let
\[
\widetilde{U}_0:=U\cap B_r(x_0),
\]
for $r>0$ sufficiently small that $\widetilde U_0$ is connected and
$0\notin \overline{\widetilde U_0}$. Since \(\Psi_U\) is locally radial
in \(U\), Lemma~\ref{lem:local-to-global-radial} applied to
$W=\widetilde U_0$ yields a scalar function \(\psi_U\) on
$I_{\widetilde U_0}$ such that
\[
\Psi_U(x)=\psi_U(|x|),\qquad x\in \widetilde{U}_0.
\]

On the other hand, by \eqref{eq:potential} and standard regularity for
logarithmic potentials with $C^2$ boundary, the gradient $\nabla\Psi_U$
extends continuously across $\partial U$, and for $z\in\partial U$ we
have
\[
\nabla \Psi_U(z)=c\,\frac{z}{|z|^2}
\]
for the same constant $c=c(\lambda,k)\in \R$ as in
\eqref{eq:potential}. Combining this with $\Psi_U(x)=\psi_U(|x|)$ on
$\widetilde U_0$ and passing to $z\in \partial \widetilde U_0\cap
\partial U$ from the inside,
\[
\psi_U'(|z|)\,\frac{z}{|z|}=c\,\frac{z}{|z|^2},
\qquad z\in \partial U\cap \overline{B_r(x_0)},
\]
hence
\[
\psi_U'(r)=\frac{c}{r}
\]
for every $r$ in the set of radii attained by points of
$\partial U\cap \overline{B_r(x_0)}$. Since $\partial U$ is a $C^2$
curve, we may, after shrinking $r$ once more, parametrize
$\partial U\cap B_r(x_0)$ by a $C^2$ embedding
$\gamma:(-\varepsilon,\varepsilon)\to \R^2$ with $\gamma(0)=x_0$.
Because $\partial U$ is not a circle centered at $0$, the function
$t\mapsto |\gamma(t)|$ is not constant on any subinterval; by choosing
$x_0\in\partial U\cap A$ at which $|\gamma|$ is non-constant on every
neighbourhood (such points form an open dense subset of $\partial
U\cap A$ when $\partial U$ is not a circle centered at $0$), we may
further assume that the derivative
$\frac{d}{dt}|\gamma(t)|\big|_{t=0}\ne 0$. Hence by the open-mapping
property for $C^1$ functions of one variable with nonzero derivative,
the image of $|\gamma|$ contains a nontrivial open interval
\(I\subset (r_1,r_2)\).
Therefore
\[
\psi_U'(r)=\frac{c}{r}
\qquad\text{for every }r\in I,
\]
whence
\[
\psi_U(r)=c\log r+C_0,
\qquad r\in I,
\]
for some real constant $C_0$.

Let
\[
W_I:=\{x\in \widetilde U_0:\ |x|\in I\}.
\]
Then \(W_I\) is a nonempty open subset of \(U\), and on \(W_I\),
\[
\Psi_U(x)=c\log|x|+C_0.
\]
Hence
\[
\Delta \Psi_U\equiv 0 \qquad \text{in } W_I,
\]
since the origin is excluded from $\widetilde U_0$. On the other hand,
distribution-theoretically,
\[
\Delta \Psi_U=2\pi\, f(|z|)\,\mathbf{1}_U-c\,\delta_0
\]
in $\C$, so on the open set $W_I\subset U\setminus\{0\}$ we have
\[
\Delta \Psi_U=2\pi\, f(|z|)=2\pi\, |z|^{2k}e^{-\pi |z|^2}.
\]
The right-hand side vanishes only at the origin, which lies outside
$W_I$, while the left-hand side just shown to be $0$ on $W_I$ — a
contradiction. This rules out the case in which $U$ is not a disk
centered at $0$, completing the proof.
\end{proof}

\begin{remark}
The structure of the argument above is in the spirit of
\cite{Aharonov-Schiffer-Zalcman}, where logarithmic-potential
rigidity statements were derived from Schiffer-style overdetermined
boundary conditions, and of the quadrature-domain literature
\cite{Gustafsson,Gustafsson-Shahgholian,Shahgholian,Karp}, in which
the regularity of free boundaries and the matching of an interior
potential with the Newtonian potential of a radially symmetric
density play a central role. Concretely, identity
\eqref{eq:potential} can be read as saying that $U$ is a
\emph{generalized quadrature domain} for the radial weight $f(|z|)$,
with quadrature data concentrated at the origin. From this perspective,
the conclusion that $U$ is a disk is a (sharp) instance of the
classical principle that radial quadrature data forces a radial
domain. The role of the $C^2$ boundary regularity is concentrated
in the boundary extension of $\nabla\Psi_U$ and in the
open-interval-of-radii argument that ends the proof; both can be
relaxed at the cost of additional technicalities, which we will
revisit in the second proof.
\end{remark}

